\documentclass[11pt]{article}
\usepackage[margin=0.9in]{geometry}
\usepackage{amsmath,amssymb,amsthm}
\usepackage{graphicx}
\usepackage{microtype}
\usepackage{xcolor}
\usepackage{hyperref}
\hypersetup{colorlinks=true,linkcolor=blue!55!black,urlcolor=blue!55!black,citecolor=blue!55!black}
\newtheorem{theorem}{Theorem}
\newtheorem{lemma}[theorem]{Lemma}
\newcommand{\Lop}{\mathcal L}
\newcommand{\B}{\mathrm B}
\title{$\tau(p;q)$-summing and quasi-$\tau(p;q)$-summing\\operators coincide isometrically}
\author{Candidate full-solution packet}
\date{13 August 2026}

\begin{document}
\maketitle

\begin{center}
\fbox{\parbox{0.92\linewidth}{\textbf{Status:} candidate full proof likely valid; human review requested.  The open problem in arXiv:1608.06926 has an affirmative isometric answer.  More strongly, the equality holds for all mixed exponents $1\le q\le p<\infty$ and in every multilinear degree.}}
\end{center}

\section{Source question and result}

Botelho and Mujica introduced quasi-$\tau(p)$-summing operators in their
duality theory for $\sigma(p)$-nuclear operators~\cite{BM17}.  At the end of
the paper they asked:
\begin{quote}
``Is it true that $\Lop_{\tau(p)}(E;F')=\Lop_{q\tau(p)}(E;F')$ for all Banach
spaces $E$ and $F$?''
\end{quote}
The question appears on page 15 of the arXiv version (page 236 of the
published version):
\begin{center}
\includegraphics[width=0.96\linewidth]{figures/open_problem_crop.png}
\end{center}

We prove an affirmative answer with equality of the defining norms.  The
argument also answers the corresponding mixed-exponent and multilinear
question, which sharpens the statement asked in the source.

\section{Definitions}

Let $E_1,\ldots,E_n,F$ be real or complex Banach spaces and
$1\le q\le p<\infty$.  For
\[
 S\in\Lop(E_1,\ldots,E_n;F')
\]
the quasi-$\tau(p;q)$ inequality is
\begin{align}
 &\left(\sum_{j=1}^m
   |S(x_{1j},\ldots,x_{nj})(y_j)|^p\right)^{1/p} \notag\\
 &\quad\le C\sup_{\substack{x_i'\in\B_{E_i'}\\y'\in\B_{F'}}}
 \left(\sum_{j=1}^m
 |x_1'(x_{1j})\cdots x_n'(x_{nj})y'(y_j)|^q\right)^{1/q},
 \label{eq:quasi}
\end{align}
for all finite families $x_{ij}\in E_i$ and $y_j\in F$.  The least $C$ is
$\|S\|_{q\tau(p;q)}$~\cite[Definition 3.1]{BM17}.

Mujica's ordinary $\tau(p;q)$ definition~\cite[Definition 3.2]{Mujica08},
specialized to a map with target $F'$, requires instead
\begin{align}
 &\left(\sum_{j=1}^m
 |y_j''(S(x_{1j},\ldots,x_{nj}))|^p\right)^{1/p} \notag\\
 &\quad\le C\sup_{\substack{x_i'\in\B_{E_i'}\\y'\in\B_{F'}}}
 \left(\sum_{j=1}^m
 |x_1'(x_{1j})\cdots x_n'(x_{nj})y_j''(y')|^q\right)^{1/q},
 \label{eq:ordinary}
\end{align}
for arbitrary $y_j''\in F''$.  The least $C$ is
$\|S\|_{\tau(p;q)}$.  Thus the quasi definition is exactly the restriction of
the ordinary definition to $y_j''=J_Fy_j$, where $J_F:F\to F''$ is the
canonical embedding.  In particular,
\begin{equation}
 \|S\|_{q\tau(p;q)}\le \|S\|_{\tau(p;q)}.
 \label{eq:easy}
\end{equation}

\section{Proof intuition}

Only a finite family of bidual test vectors occurs in any instance of
\eqref{eq:ordinary}.  Local reflexivity transfers their finite-dimensional
span into $F$ with arbitrarily small norm loss, while preserving exactly the
pairings with the finitely many values of $S$.  The ordinary numerator then
becomes a quasi numerator without error.

The denominator is subtler.  Composing a functional in $F'$ with the local
reflexivity map initially gives a functional only on the chosen subspace of
$F''$.  Hahn--Banach extends it to $F'''$, and Goldstine density approximates
that extension, on the finite test family, by canonical elements of $F'$.
This turns the operator norm bound from local reflexivity into the required
denominator bound.  The loss is $1+\varepsilon$, so taking
$\varepsilon\downarrow0$ gives an isometry.

\section{The denominator transfer}

We use the standard principle of local reflexivity in the following form:
if $M\subset F''$ and $N\subset F'$ are finite-dimensional and
$\varepsilon>0$, there is a linear $T:M\to F$ such that
\begin{equation}
 \|T\|\le1+\varepsilon,
 \qquad f'(Ty'')=y''(f')\quad(y''\in M,\ f'\in N).
 \label{eq:plr}
\end{equation}
See, for example, Oja and P\~oldvere~\cite{OP07}.

\begin{lemma}[Finite denominator transfer]\label{lem:denominator}
Let $y_1'',\ldots,y_m''\in F''$, let
$M=\operatorname{span}\{y_1'',\ldots,y_m''\}$, and let
$T:M\to F$ satisfy $\|T\|\le1+\varepsilon$.  For arbitrary
$x_{ij}\in E_i$,
\begin{align*}
 &\sup_{\substack{x_i'\in\B_{E_i'}\\f'\in\B_{F'}}}
 \left(\sum_{j=1}^m
 |x_1'(x_{1j})\cdots x_n'(x_{nj})f'(Ty_j'')|^q\right)^{1/q}\\
 &\quad\le(1+\varepsilon)
 \sup_{\substack{x_i'\in\B_{E_i'}\\g'\in\B_{F'}}}
 \left(\sum_{j=1}^m
 |x_1'(x_{1j})\cdots x_n'(x_{nj})y_j''(g')|^q\right)^{1/q}.
\end{align*}
\end{lemma}

\begin{proof}
Fix $f'\in\B_{F'}$.  The functional $f'\circ T$ on $M$ has norm at most
$1+\varepsilon$.  By Hahn--Banach it extends to a
$\Phi\in F'''$ with $\|\Phi\|\le1+\varepsilon$.  Goldstine's theorem applied
to $F'$ says that $J_{F'}(\B_{F'})$ is weak-star dense in
$\B_{F'''}$.  Consequently there is a net
$g_\alpha'\in(1+\varepsilon)\B_{F'}$ such that
\[
 y_j''(g_\alpha')\longrightarrow \Phi(y_j'')=f'(Ty_j'')
 \qquad(j=1,\ldots,m).
\]
For every fixed tuple $x_i'\in\B_{E_i'}$, continuity of a finite
$\ell_q$ norm and homogeneity now give
\begin{align*}
 &\left(\sum_{j=1}^m
 |x_1'(x_{1j})\cdots x_n'(x_{nj})f'(Ty_j'')|^q\right)^{1/q}\\
 &\quad\le(1+\varepsilon)
 \sup_{g'\in\B_{F'}}
 \left(\sum_{j=1}^m
 |x_1'(x_{1j})\cdots x_n'(x_{nj})y_j''(g')|^q\right)^{1/q}.
\end{align*}
Taking the remaining suprema proves the lemma.
\end{proof}

\section{Main theorem}

\begin{theorem}\label{thm:main}
For all real or complex Banach spaces $E_1,\ldots,E_n,F$ and all
$1\le q\le p<\infty$,
\[
 \Lop_{\tau(p;q)}(E_1,\ldots,E_n;F')
 =\Lop_{q\tau(p;q)}(E_1,\ldots,E_n;F')
\]
isometrically.  In particular, for all Banach spaces $E,F$,
\[
 \boxed{\ \Lop_{\tau(p)}(E;F')
 =\Lop_{q\tau(p)}(E;F')\quad\text{isometrically}.\ }
\]
\end{theorem}

\begin{proof}
Only the reverse of \eqref{eq:easy} needs proof.  Let $S$ be
quasi-$\tau(p;q)$-summing and fix finite families $x_{ij}\in E_i$ and
$y_j''\in F''$.  Put
\[
 M=\operatorname{span}\{y_1'',\ldots,y_m''\},\qquad
 N=\operatorname{span}\{S(x_{1j},\ldots,x_{nj}):1\le j\le m\}\subset F'.
\]
For $\varepsilon>0$, choose $T:M\to F$ as in
\eqref{eq:plr}, and set $z_j=Ty_j''$.  Preservation of the pairings gives
\[
 S(x_{1j},\ldots,x_{nj})(z_j)
 =y_j''(S(x_{1j},\ldots,x_{nj}))
 \quad(j=1,\ldots,m).
\]
Apply the quasi inequality \eqref{eq:quasi} to the families $x_{ij}$ and
$z_j$.  Lemma~\ref{lem:denominator} bounds its right-hand side by
$1+\varepsilon$ times the right-hand supremum in
\eqref{eq:ordinary}.  Hence $S$ satisfies \eqref{eq:ordinary} with constant
\[
 (1+\varepsilon)\|S\|_{q\tau(p;q)}.
\]
Since $\varepsilon>0$ is arbitrary,
$\|S\|_{\tau(p;q)}\le\|S\|_{q\tau(p;q)}$.  Together with
\eqref{eq:easy}, this proves equality of the spaces and norms.
\end{proof}

\section{Consequences and verification notes}

The source paper needed $F$ reflexive to identify the two classes.  Theorem
\ref{thm:main} removes that assumption completely.  Substituting the theorem
into Botelho--Mujica's Borel-transform result therefore replaces
quasi-$\tau(p)$ by ordinary $\tau(p)$ for arbitrary $F$, under exactly the
same bounded-approximation hypotheses on the domain duals as in their
Theorem 3.4.

The proof is valid for $p=q=1$, for all finite $p,q$ in the stated range, and
over both scalar fields.  No sequential compactness is used: Goldstine
provides a net, and only its values on finitely many vectors enter.  If the
ordinary denominator is zero, Lemma~\ref{lem:denominator} forces the
transferred quasi denominator to be zero, so that edge case is also covered.
There is no computational dependency.

Human review should concentrate on two points: (i) the specialization of
Mujica's definition to target $F'$, which makes the ordinary test vectors
members of $F''$; and (ii) Lemma~\ref{lem:denominator}, especially the use of
Goldstine for the canonical embedding $F'\hookrightarrow F'''$.  These are
the only type-sensitive steps.

\section{Bounded novelty search and limitations}

The search through 13 August 2026 included the run's result and attempt
indexes; exact-title and arXiv-id searches; the phrases
``quasi-$\tau(p)$-summing'', ``$\tau(p)$-summing local reflexivity'', and
``ideal extensions quasi-$\tau(p)$''; Botelho and Mujica's publications; and
all five works indexed by OpenAlex as citing the 2017 article.  The authors'
2019 follow-up~\cite{BM19} restates the same isometric question and explicitly
treats its answer as unknown.  No searched later source states the theorem
above or a local-reflexivity proof of it.

Accordingly this is classified as a candidate new full solution, not a known
literature answer.  The search is bounded rather than exhaustive, so novelty
confidence is moderate and specialist bibliographic review remains necessary.
The mathematical argument is self-contained modulo the standard local
reflexivity, Hahn--Banach, and Goldstine theorems.

\begin{thebibliography}{9}
\bibitem{BM17}
G. Botelho and X. Mujica,
\emph{Spaces of $\sigma(p)$-nuclear linear and multilinear operators on
Banach spaces and their duals}, Linear Algebra Appl. 519 (2017), 219--237;
arXiv:1608.06926.  DOI: \href{https://doi.org/10.1016/j.laa.2017.01.002}{10.1016/j.laa.2017.01.002}.

\bibitem{Mujica08}
X. Mujica,
\emph{$\tau(p;q)$-summing mappings and the domination theorem},
Portugaliae Math. 65 (2008), 211--226.
DOI: \href{https://doi.org/10.4171/PM/1806}{10.4171/PM/1806}.

\bibitem{OP07}
E. Oja and M. P\~oldvere,
\emph{Principle of local reflexivity revisited},
Proc. Amer. Math. Soc. 135 (2007), 1081--1088.
DOI: \href{https://doi.org/10.1090/S0002-9939-06-08612-6}{10.1090/S0002-9939-06-08612-6}.

\bibitem{BM19}
G. Botelho and X. Mujica,
\emph{Ideal extensions of classes of linear operators},
Studia Math. 247 (2019), 285--297.
DOI: \href{https://doi.org/10.4064/sm171030-15-3}{10.4064/sm171030-15-3}.
\end{thebibliography}

\end{document}
